\chapter{Problem Definition and Data Protocol}
\label{chap:problem_data}
\updateheader{Chapter 3. Problem Definition and Data Protocol}

%==============================================================================
% 本章作用：定义任务（输入/输出/约束），并说明评测数据如何得到、为什么可信。
% 写作依据：docs/DATA_AUDIT.md, docs/DATA_PREPARATION.md,
%           data/versions/thesis-v0.1/{README.md, manifest.json}
% 所有数字均来自 thesis-v0.1 和 thesis-v0.2 的 manifest.json。
%==============================================================================

This chapter defines the question answering task studied in this thesis and describes the data protocol used for development and evaluation. Section~\ref{sec:task_definition} specifies the inputs, outputs, and constraints of the task. Section~\ref{sec:data_source} introduces the source data. Section~\ref{sec:inherited_audit} reports an audit of the evaluation material inherited from the earlier conference prototype and explains why that material is not used directly as a thesis test set. Section~\ref{sec:dataset_construction} describes the construction of the versioned dataset \dataversion{}, and Section~\ref{sec:review_protocol} defines the review and freezing procedure for the final evaluation set. Section~\ref{sec:leakage_control} summarizes the measures taken against information leakage.

The data protocol is presented before the proposed method for two reasons. First, the evaluation claims of this thesis depend on a clear separation between development questions, evaluation questions, and reference answers. Second, the audit revealed several properties of the inherited benchmark, such as cross-split duplicates and answer-aware candidate selection, that would otherwise weaken the interpretation of later experiments.

%==============================================================================
\section{Task Definition}
\label{sec:task_definition}

The task is multiple-choice cardiology question answering with an explicit, evidence-linked explanation. Each instance consists of a clinical question stem $q$ and a set of four answer options $\mathcal{O} = \{o_A, o_B, o_C, o_D\}$. The system may consult a fixed evidence corpus $\mathcal{C}$ but has no access to the reference answer, to relevance labels, or to any explanation written for the question.

\begin{definition}[Evidence-grounded cardiology question answering]
\label{def:task}
Given a question $q$, its options $\mathcal{O}$, and an evidence corpus $\mathcal{C}$, a system $f$ returns
\begin{equation}
  f(q, \mathcal{O}, \mathcal{C}) = (\hat{a},\, e,\, G,\, \tau),
  \label{eq:task_output}
\end{equation}
where $\hat{a} \in \mathcal{O} \cup \{\varnothing\}$ is the selected option or an explicit abstention, $e$ is a natural-language explanation, $G$ is an OPM-style graph describing the mechanism that connects the question findings to the selected answer, and $\tau$ is an execution trace recording model calls, retrieved passages, verification findings, and repair attempts.
\end{definition}

The graph $G$ is defined over objects, processes, and states, with typed relations between them. Each element and relation may carry references to passages in $\mathcal{C}$ that support it. The precise representation is given in Chapter~\ref{chap:method}. In this chapter, it is sufficient to note that the task has two separable aspects: \emph{answer correctness}, which is judged against the reference option, and \emph{explanation support}, which is judged by whether the graph and explanation are supported by the cited evidence and consistent with the selected answer.

Three constraints follow from this definition.
\begin{enumerate}[label=(\roman*)]
  \item \textbf{Answer isolation.} Reference answers, answer indices, and explanations must never enter the system input or the retrieval corpus.
  \item \textbf{Fixed evidence.} For comparisons between system variants, the evidence corpus and retrieval configuration are fixed in advance, so that differences in output cannot be attributed to different evidence.
  \item \textbf{Explicit failure.} Abstentions, invalid outputs, and missing predictions are counted explicitly rather than removed from the evaluation.
\end{enumerate}

The task does not include clinical decision-making. The outputs of the system are research artifacts intended for auditing explanation behavior; they are not diagnostic or treatment recommendations.

%==============================================================================
\section{Source Data}
\label{sec:data_source}

The questions are drawn from MedQA, a multiple-choice question answering dataset constructed from medical licensing examinations~\cite{jin_medqa_2021}. This thesis uses the English questions of the United States Medical Licensing Examination (USMLE) style, each with four answer options.

The local copy of the raw data contains 11,451 distinct question--option records, of which 10,178 are labeled as training questions and 1,273 as test questions by the source. These counts correspond to the training and test partitions of the English MedQA subset; no separate development file was present in the local copy. The checksum of the raw file is recorded in the dataset manifest so that every derived file can be traced back to the same source. The provenance of the upstream release itself is not independently certified in this thesis.

Two properties of the raw data are relevant to later sections. First, the raw file contains two groups of records that share an identical question stem but differ in their options. Second, no identical question stem appears in both the source training and the source test partition. Table~\ref{tab:raw_stats} summarizes these statistics.

\begin{table}[H]
\centering
\caption{Statistics of the local MedQA source file.}
\label{tab:raw_stats}
\begin{tabular}{lr}
\toprule
Item & Count \\
\midrule
Distinct question--option records & 11,451 \\
\quad Source-reported training questions & 10,178 \\
\quad Source-reported test questions & 1,273 \\
Options per question & 4 \\
Stem groups with more than one option set & 2 \\
Stem groups spanning both source partitions & 0 \\
\bottomrule
\end{tabular}
\end{table}

%==============================================================================
\section{Audit of the Inherited Evaluation Material}
\label{sec:inherited_audit}

The thesis builds on an earlier conference prototype that produced OPM-style graphs for cardiology questions. That project left three kinds of material: processed candidate lists of cardiology-related questions, manual relevance annotations, and a packaged benchmark with a development split and a test split of 100 questions each. Before reusing any of this material, it was audited against the raw source. The audit did not modify any inherited file or label.

\subsection{Candidate Lists}

The two inherited candidate lists contain 546 and 935 records, respectively, giving 1,481 occurrences in total. Every record in the smaller list also appears in the larger one, so the union contains 935 distinct questions rather than 1,481. Each record was matched to exactly one raw record, and its source partition and answer key agree with the raw data.

\subsection{Packaged Benchmark}

The packaged benchmark has internally unique identifiers and consistent labels, but three issues limit its use as an independent test set.

\begin{enumerate}[label=(\alph*)]
  \item \textbf{Mixed source partitions.} The packaged development split contains 93 questions from the source training partition and 7 from the source test partition. The packaged test split contains 88 and 12, respectively. The package names therefore do not correspond to the original MedQA partitions.
  \item \textbf{Cross-split duplicates.} Twelve questions appear in both the packaged development and test splits with identical stems and identical options.
  \item \textbf{Conflicting labels.} Two of these duplicated questions received opposite cardiology-relevance labels in the two splits. In addition, 35 identifier strings occur in both splits. Because the identifiers were assigned separately for each candidate list, an identical identifier string does not indicate the same question, and records must be linked by the pair (split, identifier) or by content.
\end{enumerate}

\subsection{Manual Annotations}

The 200 inherited annotation rows cover 188 distinct questions. Of these, 134 received consistently positive cardiology-relevance judgments, 52 consistently negative judgments, and 2 conflicting judgments. The records do not document reviewer qualifications or inter-rater agreement. These annotations are therefore preserved and used as inherited labels, not as expert adjudications.

\subsection{Answer-Aware Candidate Selection}

The keyword filter used to build the inherited candidate lists reads not only the question but also the reference answer, the answer index, the explanation, and the options. A question may therefore have been selected partly because its \emph{answer} mentions a cardiological concept. Such selection is acceptable for building a topic-specific pool, but it is not a blind evaluation protocol: it can favor questions whose answers align with cardiology keywords. Because it cannot be established which code revision produced each historical artifact, the inherited candidate lists are used only for development in this thesis.

Table~\ref{tab:audit_summary} summarizes the findings and their consequences.

\begin{table}[H]
\centering
\caption{Audit findings for the inherited material and their consequences.}
\label{tab:audit_summary}
\begin{tabularx}{\linewidth}{>{\raggedright\arraybackslash}X >{\raggedright\arraybackslash}X}
\toprule
Finding & Consequence in this thesis \\
\midrule
1,481 candidate occurrences correspond to 935 distinct questions & Deduplicate by content-based identity before any split \\
Packaged splits mix source training and test questions & Do not treat package names as dataset partitions \\
12 questions duplicated across packaged splits & Packaged test split is not an independent test set \\
2 conflicting relevance labels & Quarantine; adjudicate only with documented evidence \\
Candidate selection read answer and explanation fields & Use inherited candidates for development only; build a new question-only evaluation pool \\
\bottomrule
\end{tabularx}
\end{table}

%==============================================================================
\section{Construction of the Thesis Dataset}
\label{sec:dataset_construction}

Based on the audit, a new dataset version, \dataversion{}, was constructed with a deterministic script. Figure~\ref{fig:data_pipeline} gives an overview of the construction process.

\begin{figure}[H]
\centering
\begin{tikzpicture}[
  node distance=5mm and 8mm,
  box/.style={draw, rounded corners=2pt, align=center, font=\small, minimum width=40mm, minimum height=9mm, fill=gray!8},
  key/.style={box, fill=blue!8},
  arr/.style={-{Stealth[length=2mm]}, thick}
]
\node[box] (raw) {Local MedQA source\\11,451 records};

\node[box, below left=8mm and -12mm of raw] (legacy) {Inherited candidates\\935 distinct questions};
\node[box, below=of legacy] (ann) {Inherited annotations\\188 questions};
\node[box, below=of ann] (elig) {Consistent positive,\\training-only, no variants\\118 questions};
\node[key, below=of elig] (pilot) {Development pilot\\20 questions};

\node[box, below right=8mm and -12mm of raw] (screen) {Question-only screening\\of unexposed stems};
\node[box, below=of screen] (cand) {Test candidates: 51\\(training candidates: 456)};
\node[box, below=of cand] (review) {Manual review\\scope, visuals, duplicates};
\node[key, below=of review] (test) {Frozen evaluation set\\28 questions};

\draw[arr] (raw) -- (legacy);
\draw[arr] (legacy) -- (ann);
\draw[arr] (ann) -- (elig);
\draw[arr] (elig) -- (pilot);
\draw[arr] (raw) -- (screen);
\draw[arr] (screen) -- (cand);
\draw[arr] (cand) -- (review);
\draw[arr] (review) -- (test);
\draw[arr, dashed] (legacy.east) -- node[above, font=\scriptsize, align=center]{excluded\\stems} (screen.west);
\end{tikzpicture}
\caption{Construction of the development pilot (left) and the prospective evaluation pool (right) in \dataversion{}.}
\label{fig:data_pipeline}
\end{figure}

\subsection{Content-Based Identifiers}

The identifiers of the inherited material are local to each candidate list and cannot be used to link questions across files. Every question is therefore assigned a canonical identifier derived from its content:
\begin{equation}
  \mathrm{id}(q, \mathcal{O}) = \texttt{mq\_} \,\Vert\, \mathrm{SHA256}\big(\mathrm{norm}(q),\, \mathrm{sort}(\mathrm{norm}(\mathcal{O}))\big),
  \label{eq:canonical_id}
\end{equation}
where $\mathrm{norm}(\cdot)$ applies Unicode NFC normalization and collapses whitespace while preserving case, punctuation, and numbers, and the options are sorted by their keys. Two records receive the same identifier only if their normalized stems and options are identical. A separate stem identifier, computed from the normalized stem alone, is used to detect records that share a stem but differ in their options.

\subsection{Development Pilot}

The development pilot is used to build and debug the system and to tune prompts. It is selected only from questions whose inherited annotations are consistent and positive and whose source partition is training. Among the 167 reviewed training-only questions with consistent labels, 119 are positive. One of them shares its stem with another raw record that has different options; it is excluded because the variant has not been resolved. This leaves 118 eligible questions.

Twenty questions are selected deterministically from this pool: each eligible question is ranked by the SHA-256 digest of the pair (seed, identifier) with seed 42, and the first 20 are taken. This procedure is reproducible and does not depend on any answer, label, or model output. Because the pilot is used during development, results on it are never reported as held-out performance.

\subsection{Prospective Evaluation Pool}

The evaluation pool is built with a new screening procedure that reads the question stem only. A question is a candidate if its stem contains at least one of 30 cardiology terms (Table~\ref{tab:screening_terms}), matched as whole words after normalization and case folding. The term list is recorded in the dataset manifest. Options, answers, and explanations are never read during screening.

\begin{table}[H]
\centering
\caption{Terms used for question-only candidate screening.}
\label{tab:screening_terms}
\small
\begin{tabularx}{\linewidth}{>{\raggedright\arraybackslash}X}
\toprule
angina; aortic coarctation; aortic stenosis; arrhythmia; atrial fibrillation; cardiac arrest; cardiomyopathy; coarctation; coarctation of the aorta; coronary artery disease; ECG; echocardiography; EKG; endocarditis; heart failure; mitral prolapse; mitral regurgitation; mitral valve prolapse; murmur; myocardial infarction; myocarditis; patent ductus arteriosus; PDA; pericarditis; pulmonary embolism; pulmonary embolus; tetralogy of Fallot; TOF; valvular disease; ventricular tachycardia \\
\bottomrule
\end{tabularx}
\end{table}

Two exclusion rules are applied before screening. First, every stem that appeared in the inherited candidate lists or in earlier graph annotations is excluded, including stems that appear with different options. This ensures that no question seen during the conference work, or during its manual analysis, enters the evaluation pool. Second, the source partition of each record is preserved.

Screening produced 51 candidates from the source test partition and 456 from the source training partition. The 51 test-partition candidates form the primary evaluation pool. The training-partition candidates are counted but not labeled; they are a reserve that can be reviewed under the same protocol if the primary pool proves too small. The reserve was not used for the frozen evaluation set reported in this thesis.

Keyword membership is not a judgment of clinical relevance. A question may mention a murmur while testing an unrelated concept, and a cardiology question may use none of the listed terms. The candidates therefore require manual review (Section~\ref{sec:review_protocol}).

\subsection{Near-Duplicate and Visual-Material Flags}

Exact identifiers do not detect questions that differ only slightly in wording. Each test candidate is compared with every other question in the inherited candidate pool and in the new candidate lists (pairs with an identical stem are handled by the stem identifier instead), and a pair is flagged as a possible near duplicate if it satisfies all of the following: the ratio of the shorter to the longer normalized stem length is at least 0.85, the Jaccard similarity of their word sets is at least 0.90, and their character-level sequence similarity is at least 0.92. One pair was flagged. It spans the source test and training partitions and requires manual inspection. These lexical thresholds generate review candidates only; they cannot establish that all semantic duplicates have been found.

In addition, 9 of the 51 test candidates mention a figure, image, photograph, diagram, or radiograph, or refer to material ``shown'' above or below. Such questions may be unanswerable if the referenced material is missing. The flag prompts inspection; it does not by itself make a question invalid. Manual review showed that the flag is imperfect in both directions: in four flagged questions the referenced material was given in the text or was not needed, while two unflagged questions referred to missing material with wording that the pattern does not match (``picture'' and the plural ``Figures''). Exhibit dependence was therefore judged for every candidate during review rather than taken from the flag.

\subsection{Released Files}

Table~\ref{tab:dataset_files} lists the files of \dataversion{} and their intended use. Model-input files and evaluator-only files are stored separately.

\begin{table}[H]
\centering
\caption{Files of the dataset version \dataversion{} and their intended use.}
\label{tab:dataset_files}
\footnotesize
\begin{tabularx}{\linewidth}{>{\raggedright\arraybackslash}p{0.50\linewidth} >{\raggedright\arraybackslash}X}
\toprule
File & Intended use \\
\midrule
\texttt{inputs/pilot\_questions.jsonl} & System input for development (20 questions) \\
\texttt{evaluator\_only/pilot\_gold.jsonl} & Scoring of pilot outputs; not reported as held-out \\
\texttt{inputs/reviewed\_train\_questions.jsonl} & Development pool with positive and negative examples (167 questions) \\
\texttt{inputs/prospective\_test\_questions.jsonl} & Evaluation candidates (51 questions, before review) \\
\texttt{review/prospective\_test\_screening.csv} & Review form for the evaluation candidates \\
\texttt{review/relevance\_conflicts.csv} & Review form for the two conflicting labels \\
\texttt{review/near\_duplicate\_pairs.csv} & Side-by-side comparison of the flagged pair, without answer keys \\
\texttt{manifest.json}, \texttt{MANIFEST.sha256} & Selection policy, counts, and file checksums \\
\bottomrule
\end{tabularx}
\end{table}

%==============================================================================
\section{Review and Freezing Protocol}
\label{sec:review_protocol}

\subsection{Scope Definition}

A question is in scope if its \emph{main tested concept} is cardiovascular. Table~\ref{tab:scope_rules} lists the final scope rules, including the decisions for boundary cases. The rules were fixed by the author before the final labels were assigned. Several boundary rules were added after reading the candidate questions; they were not informed by answer keys or by any system output.

\begin{table}[H]
\centering
\caption{Final scope rules for the evaluation set.}
\label{tab:scope_rules}
\small
\begin{tabularx}{\linewidth}{>{\raggedright\arraybackslash}p{0.24\linewidth} >{\raggedright\arraybackslash}X}
\toprule
Case & Rule \\
\midrule
Core rule & A question is in scope if its main tested concept is cardiovascular. A cardiac history, vital-sign value, risk factor, or cardiovascular distractor option is not sufficient on its own. \\
Cardiovascular pharmacology & Include drug selection, treatment mechanisms, contraindications and adverse effects when the question directly tests cardiovascular care or physiology. Exclude primarily noncardiovascular drug effects and toxicology. \\
Calculations & Include if the calculation directly computes a cardiovascular physiological parameter (e.g., cardiac output, ejection fraction, systemic vascular resistance, mean arterial pressure). \\
Vascular disease & Include aortic, peripheral arterial and venous thrombotic disease, and questions that directly test a circulatory or hemodynamic mechanism. \\
Pulmonary embolism & Include pulmonary embolism and pulmonary vascular disease when they are the primary target of diagnosis, mechanism or management. \\
Secondary neurological injury & Include if the primary tested concept is cardiovascular prevention or management of the cardiac source (e.g., anticoagulation for atrial fibrillation to prevent cardioembolic stroke); exclude pure neurological localization. \\
Congenital heart disease in syndromes & Include diagnosis, embryology or etiology of a cardiac lesion. Exclude syndrome recognition and questions asking for a noncardiac associated finding. \\
Other boundary cases & Assess scope and test eligibility separately. Keep exhibits described in sufficient text. Exclude essential missing exhibits and duplicates. Use \texttt{unclear} for unresolved scope, exhibit dependence or source quality. \\
\bottomrule
\end{tabularx}
\end{table}

\todo{Rule for cardiovascular risk-factor management (one included question tests it): add an explicit clause, or explain here why the core rule covers it.}

\subsection{Review Fields}

Each candidate is reviewed with the question and options only; answer keys are not shown in the review form. For every candidate, the fields in Table~\ref{tab:review_fields} are recorded. The recorded reasons describe the tested concept only and do not state the answer.

\begin{table}[H]
\centering
\caption{Fields recorded during review of each evaluation candidate.}
\label{tab:review_fields}
\begin{tabularx}{\linewidth}{l >{\raggedright\arraybackslash}X}
\toprule
Field & Content \\
\midrule
\texttt{is\_cardiology} & \texttt{yes} or \texttt{no} under the scope rules (\texttt{unclear} allowed during drafting only) \\
\texttt{include\_for\_evaluation} & \texttt{yes} or \texttt{no}, also accounting for missing exhibits, duplicates and source quality \\
\texttt{exclusion\_reason} & Out of scope, essential missing exhibit, near duplicate across source partitions, or source quality \\
\texttt{rule\_applied} & The rule in Table~\ref{tab:scope_rules} on which the decision rests \\
\texttt{reason\_and\_evidence} & Short justification naming the tested concept \\
\texttt{reviewer}, \texttt{review\_date} & Who made the final decision and when \\
\bottomrule
\end{tabularx}
\end{table}

For the flagged near-duplicate pair, stems, numerical values, final questions, options, and source partitions are compared. For the two conflicting relevance labels, an adjudicated label is recorded together with the reviewer, date, and supporting reason; the original labels remain preserved.

\subsection{Labeling Procedure}

Candidate labels were drafted with the assistance of large language models. Three LLMs (ChatGPT, Gemini, and Claude) independently drafted labels for all 51 candidates. Each draft was produced under a different preliminary version of the scope rules. On \texttt{include\_for\_evaluation}, the three drafts agreed on 36 of 51 candidates (70.6\%); pairwise agreement ranged from 39 to 41 of 51. Because the drafts used different rule versions, these figures describe the stability of the drafting step, not inter-annotator agreement between human reviewers.

The author then fixed the final rules in Table~\ref{tab:scope_rules}, all candidates were re-labeled under these rules, and the author confirmed every final decision, including all items on which the earlier drafts had disagreed. The author is a graduate student in information science without clinical training. \todo{Confirm this description of your background.} No independent clinical expert reviewed the labels; this limitation is discussed in Chapter~\ref{chap:discussion}.

\subsection{Review Outcome and Freezing}

Table~\ref{tab:review_outcome} summarizes the outcome. Of the 51 candidates, 32 were judged in scope and 28 were included in the evaluation set. Two in-scope questions were excluded because a missing exhibit was essential for answering, one because it is a near duplicate of a source-training question, and one because its option set does not contain the histology expected for the described lesion, which raised doubts about the source item. This last decision was made without consulting the answer key.

\begin{table}[H]
\centering
\caption{Outcome of the evaluation-candidate review.}
\label{tab:review_outcome}
\begin{tabular}{lr}
\toprule
Item & Count \\
\midrule
Candidates reviewed & 51 \\
\quad Out of scope & 19 \\
\quad In scope & 32 \\
\qquad Excluded: essential missing exhibit & 2 \\
\qquad Excluded: near duplicate across source partitions & 1 \\
\qquad Excluded: source quality & 1 \\
\textbf{Included in the evaluation set} & \textbf{28} \\
\midrule
Conflicting relevance labels adjudicated (both not in scope) & 2 \\
Near-duplicate pairs resolved (same underlying item) & 1 \\
\bottomrule
\end{tabular}
\end{table}

The flagged near-duplicate pair was judged to be the same underlying item: the stems, findings and final questions are nearly identical and only the option sets differ. The source-test record is excluded, and both records are listed in a quarantine file so that neither can enter an evaluation set later, for example when the training reserve is reviewed.

The frozen evaluation set is released as dataset version \texttt{thesis-v0.2}, built deterministically from \dataversion{} and the recorded decisions. It contains the 28 evaluation questions as model inputs, their reference answers in a separate evaluator-only file, every review decision with its reason, the final scope rules, and checksums of all files. The development pilot and the development pool are carried over unchanged. The evidence corpus, the prompts, and the evaluation rules are frozen before any final experiment is run, and selection never depended on whether a system answered a question correctly.

%==============================================================================
\section{Leakage Control}
\label{sec:leakage_control}

The following measures prevent reference information from reaching the system.

\begin{itemize}
  \item \textbf{Minimal input records.} Every model-input record contains exactly three fields: the identifier, the question, and the options. Answers, answer indices, relevance labels, screening terms, prior model judgments, and annotation reasons are excluded.
  \item \textbf{Separated reference files.} Reference answers are stored in an evaluator-only directory. Because a directory name is not an access-control boundary, the experiment runner loads only the designated input file, and evaluator files are excluded from the retrieval corpus.
  \item \textbf{Exposure-aware screening.} Stems seen during the conference work are excluded from the evaluation pool, as described in Section~\ref{sec:dataset_construction}.
  \item \textbf{Duplicate quarantine.} Records judged to be the same underlying item are listed in a quarantine file and can never enter an evaluation set.
  \item \textbf{Traceability.} Each derived file is linked to the raw source by checksums and source row indices recorded in the manifest.
\end{itemize}

These measures prevent leakage at the level of this project. They cannot rule out that public MedQA questions were part of the pretraining data of the language model used in the experiments. This possibility is discussed as a threat to validity in Chapter~\ref{chap:discussion}.

%==============================================================================
\section{Chapter Summary}
\label{sec:ch3_summary}

This chapter defined the task as multiple-choice cardiology question answering with an evidence-linked OPM explanation, and separated answer correctness from explanation support. An audit of the inherited material showed that its packaged benchmark mixes source partitions, contains cross-split duplicates and conflicting labels, and was selected with access to answers. The thesis therefore constructs a new dataset version with content-based identifiers, a reproducible 20-question development pilot, and a question-only evaluation pool of 51 candidates. After LLM-assisted drafting and author adjudication under written scope rules, 28 questions form the frozen evaluation set \texttt{thesis-v0.2}. The next chapter describes the method evaluated on this data.
